\documentclass[12pt]{article}

\usepackage[margin=1in]{geometry}
\usepackage{amsmath,amssymb}
\usepackage{booktabs}
\usepackage{graphicx}
\usepackage{natbib}
\usepackage{setspace}
\usepackage{hyperref}
\usepackage[T1]{fontenc}
\usepackage{lmodern}
\usepackage{threeparttable}
\usepackage{xcolor}
\usepackage{xr-hyper}
\externaldocument[M-][nocite]{main}

\hypersetup{
    colorlinks=true,
    linkcolor=blue!60!black,
    citecolor=blue!60!black,
    urlcolor=blue!60!black,
}

\onehalfspacing

% Auto-generated numeric macros shared with main.tex.
% !TEX root = main.tex
% ------------------------------------------------------------------
% AUTO-GENERATED by scripts/export_manuscript_numbers.jl
% Do not edit by hand. Regenerate after any analysis re-run.
% ------------------------------------------------------------------

\newcommand{\pGamma}{2.5}
\newcommand{\pBeta}{0.97}
\newcommand{\pRRate}{2\%}
\newcommand{\pRRateNum}{0.02}
\newcommand{\pInflation}{2\%}
\newcommand{\pInflationNum}{0.02}
\newcommand{\pCFloor}{\$6{,}180}
\newcommand{\pFixedCost}{\$2{,}500}
\newcommand{\pMinWealth}{\$5{,}000}
\newcommand{\pWMax}{\$3{,}000{,}000}
\newcommand{\pWMaxMillions}{\$3\text{ million}}
\newcommand{\pThetaDFJ}{56.96}
\newcommand{\pKappaDFJ}{\$272{,}628}
\newcommand{\pDeltaC}{0.02}
\newcommand{\pHealthUtilGood}{1.00}
\newcommand{\pHealthUtilFair}{0.92}
\newcommand{\pHealthUtilPoor}{0.82}
\newcommand{\pHazardGood}{0.50}
\newcommand{\pHazardFair}{1.0}
\newcommand{\pHazardPoor}{3.8}
\newcommand{\pPessimism}{0.960}
\newcommand{\pAgeStart}{65}
\newcommand{\pAgeEnd}{110}
\newcommand{\pT}{46}
\newcommand{\pNWealth}{80}
\newcommand{\pNAnnuity}{30}
\newcommand{\pNAlpha}{101}
\newcommand{\pNQuad}{9}
\newcommand{\pMwrBaseline}{0.87}
\newcommand{\pMwrLoad}{13\%}
\newcommand{\nHRSTotal}{5,303}
\newcommand{\nHRSEligible}{2846}
\newcommand{\pctHRSAboveWmax}{0.4\%}
\newcommand{\nHRSLifetimeEligible}{2,860}
\newcommand{\nHRSLifetimeOwners}{89}
\newcommand{\pctHRSLifetime}{3.11\%}
\newcommand{\pctHRSIannPooled}{5.21\%}
\newcommand{\pctHRSLifetimeCILow}{2.54\%}
\newcommand{\pctHRSLifetimeCIHigh}{3.81\%}
\newcommand{\pctHRSIannCILow}{4.45\%}
\newcommand{\pctHRSIannCIHigh}{6.09\%}
\newcommand{\pctELSAWaveSixAnnuity}{90.2\%}
\newcommand{\nELSAPostLumpSum}{869}
\newcommand{\pctELSAPostAnnuity}{1.27\%}
\newcommand{\ownFrictionless}{41.8\%}
\newcommand{\ownAddSS}{100.0\%}
\newcommand{\ownAddBequests}{100.0\%}
\newcommand{\ownAddMedRS}{87.7\%}
\newcommand{\ownAddPessimism}{62.4\%}
\newcommand{\ownAddLoads}{18.2\%}
\newcommand{\ownSixChannel}{8.5\%}
\newcommand{\ownSevenChannelExt}{2.6\%}
\newcommand{\ownEightChannelExt}{1.8\%}
\newcommand{\ownNineChannelLTC}{2.0\%}
\newcommand{\ownNoBehavioralBaseline}{2.0\%}
\newcommand{\ownTenChannelSDU}{51.1\%}
\newcommand{\ownElevenChannelFull}{0.0\%}
\newcommand{\ownModelOne}{0.0\%}
\newcommand{\ownModelOneStructural}{0.0\%}
\newcommand{\pWedgeFactorLow}{0.137}
\newcommand{\pWedgeFactorMid}{0.179}
\newcommand{\pWedgeFactorHigh}{0.263}
\newcommand{\pUKRetentionPre}{95\%}
\newcommand{\pUKRetentionLow}{13\%}
\newcommand{\pUKRetentionMid}{17\%}
\newcommand{\pUKRetentionHigh}{25\%}
\newcommand{\ownFrictionlessNum}{41.8}
\newcommand{\ownWedgeLow}{5.7\%}
\newcommand{\ownWedgeMid}{7.5\%}
\newcommand{\ownWedgeHigh}{11.0\%}
\newcommand{\ownModelTwoLow}{5.7\%}
\newcommand{\ownModelTwoMid}{7.5\%}
\newcommand{\ownModelTwoHigh}{11.0\%}
\newcommand{\ownModelTwo}{7.5\%}
\newcommand{\ownHeadline}{7.5\%}
\newcommand{\ownBracketLow}{5.7\%}
\newcommand{\ownBracketHigh}{11.0\%}
\newcommand{\retentionSS}{239.0\%}
\newcommand{\ownPrePessimism}{87.7\%}
\newcommand{\deltaPessimism}{-25.3}
\newcommand{\retentionPessimism}{71.1\%}
\newcommand{\retentionLoads}{29.1\%}
\newcommand{\retentionInflation}{46.6\%}
\newcommand{\deltaMedRS}{-12.3}
\newcommand{\magDeltaMedRS}{12.3}
\newcommand{\retentionMedRS}{87.7\%}
\newcommand{\deltaAgeNeeds}{-5.8}
\newcommand{\deltaStateUtil}{-0.9}
\newcommand{\shapSS}{-11.5}
\newcommand{\shapBequests}{0.6}
\newcommand{\shapMedRS}{15.4}
\newcommand{\shapPessimism}{5.4}
\newcommand{\shapAgeNeeds}{2.5}
\newcommand{\shapStateUtil}{0.3}
\newcommand{\shapLoads}{16.2}
\newcommand{\shapInflation}{1.2}
\newcommand{\shapLTC}{0.2}
\newcommand{\shapSDU}{-30.4}
\newcommand{\shapPED}{41.9}
\newcommand{\shapNarrowFraming}{41.9}
\newcommand{\shapShareSS}{-27\%}
\newcommand{\shapShareBequests}{1\%}
\newcommand{\shapShareMedRS}{37\%}
\newcommand{\shapSharePessimism}{13\%}
\newcommand{\shapShareAgeNeeds}{6\%}
\newcommand{\shapShareStateUtil}{1\%}
\newcommand{\shapShareLoads}{39\%}
\newcommand{\shapShareInflation}{3\%}
\newcommand{\shapShareLTC}{1\%}
\newcommand{\shapShareSDU}{-73\%}
\newcommand{\shapSharePED}{100\%}
\newcommand{\shapRSPessAge}{23.3}
\newcommand{\shapShareRSPessAge}{56\%}
\newcommand{\shapNineSS}{-24.8}
\newcommand{\shapShareNineSS}{-62\%}
\newcommand{\shapNineBequests}{5.4}
\newcommand{\shapShareNineBequests}{14\%}
\newcommand{\shapNineMedRS}{11.5}
\newcommand{\shapShareNineMedRS}{29\%}
\newcommand{\shapNinePessimism}{11.6}
\newcommand{\shapShareNinePessimism}{29\%}
\newcommand{\shapNineAgeNeeds}{4.6}
\newcommand{\shapShareNineAgeNeeds}{12\%}
\newcommand{\shapNineStateUtil}{0.2}
\newcommand{\shapShareNineStateUtil}{1\%}
\newcommand{\shapNineLoads}{26.8}
\newcommand{\shapShareNineLoads}{67\%}
\newcommand{\shapNineInflation}{4.4}
\newcommand{\shapShareNineInflation}{11\%}
\newcommand{\shapNineLTC}{0.1}
\newcommand{\shapShareNineLTC}{0\%}
\newcommand{\cevGoodHundredKNoBq}{0.0\%}
\newcommand{\cevGoodHundredKDFJ}{0.0\%}
\newcommand{\cevGoodHundredKStrong}{0.0\%}
\newcommand{\cevGoodTwoHundredKNoBq}{0.0\%}
\newcommand{\cevGoodTwoHundredKDFJ}{0.0\%}
\newcommand{\cevGoodTwoHundredKStrong}{0.0\%}
\newcommand{\cevGoodOneMillNoBq}{6.6\%}
\newcommand{\cevGoodOneMillDFJ}{0.8\%}
\newcommand{\cevGoodOneMillStrong}{0.0\%}
\newcommand{\popMeanCevNoBq}{0.5\%}
\newcommand{\popMeanCevDFJ}{0.0\%}
\newcommand{\popFracPosNoBq}{15.8\%}
\newcommand{\popFracPosDFJ}{7.2\%}
\newcommand{\popFracAboveOneNoBq}{11.8\%}
\newcommand{\popFracAboveOneDFJ}{0.0\%}
\newcommand{\ownGroupPricing}{6.1\%}
\newcommand{\ownPublicOption}{14.3\%}
\newcommand{\ownActuariallyFair}{25.6\%}
\newcommand{\ownRealTIPS}{0.0\%}
\newcommand{\ownRealNomEquiv}{8.1\%}
\newcommand{\ownFairReal}{31.8\%}
\newcommand{\ownCorrectPessimism}{18.3\%}
\newcommand{\ownGroupPlusCorrect}{25.9\%}
\newcommand{\ownBestFeasible}{35.0\%}
\newcommand{\cevGoodHundredKBaseline}{0.0\%}
\newcommand{\cevGoodHundredKGroupPrice}{0.0\%}
\newcommand{\cevGoodHundredKRealAnn}{0.0\%}
\newcommand{\cevGoodHundredKBestFeas}{0.0\%}
\newcommand{\cevGoodTwoHundredKBaseline}{0.0\%}
\newcommand{\cevGoodTwoHundredKGroupPrice}{0.0\%}
\newcommand{\cevGoodTwoHundredKRealAnn}{0.0\%}
\newcommand{\cevGoodTwoHundredKBestFeas}{0.7\%}
\newcommand{\cevGoodFiveHundredKBaseline}{0.0\%}
\newcommand{\cevGoodFiveHundredKGroupPrice}{0.5\%}
\newcommand{\cevGoodFiveHundredKRealAnn}{0.1\%}
\newcommand{\cevGoodFiveHundredKBestFeas}{4.1\%}
\newcommand{\cevGoodOneMillBaseline}{0.8\%}
\newcommand{\cevGoodOneMillGroupPrice}{1.5\%}
\newcommand{\cevGoodOneMillRealAnn}{1.0\%}
\newcommand{\cevGoodOneMillBestFeas}{5.6\%}
\newcommand{\ownGammaTwo}{0.0\%}
\newcommand{\ownGammaTwoPointThree}{0.0\%}
\newcommand{\ownGammaTwoPointFour}{0.1\%}
\newcommand{\ownGammaTwoPointFive}{1.8\%}
\newcommand{\ownGammaThree}{18.1\%}
\newcommand{\ownInflationOne}{4.7\%}
\newcommand{\ownInflationTwo}{1.8\%}
\newcommand{\ownInflationThree}{0.7\%}
\newcommand{\ownPsiObjective}{21.4\%}
\newcommand{\ownPsiNinetySeven}{5.1\%}
\newcommand{\ownPsiBaseline}{8.7\%}
\newcommand{\ownPsiNinetyNine}{14.0\%}
\newcommand{\ownMWREightyTwo}{0.0\%}
\newcommand{\ownMWREightyFive}{0.4\%}
\newcommand{\ownMWRNinety}{5.8\%}
\newcommand{\ownMWRNinetyFive}{15.4\%}
\newcommand{\ownBequestNone}{21.0\%}
\newcommand{\ownHazardRS}{3.2\%}
\newcommand{\ownHazardHRS}{6.4\%}
\newcommand{\ownHazardConservative}{10.6\%}
\newcommand{\ownHazardAgeBand}{7.3\%}
\newcommand{\ownSSCutZero}{1.8\%}
\newcommand{\ownSSCutTen}{3.8\%}
\newcommand{\ownSSCutFifteen}{5.9\%}
\newcommand{\ownSSCutTwentyThree}{10.4\%}
\newcommand{\ownSSCutThirty}{16.3\%}
\newcommand{\ownSSCutForty}{25.1\%}
\newcommand{\ownSSCutFifty}{28.6\%}
\newcommand{\ownSSCutHundred}{14.3\%}
\newcommand{\ownNineChannelFLN}{1.5\%}
\newcommand{\pHealthUtilFairRS}{0.95}
\newcommand{\pHealthUtilPoorRS}{0.85}
\newcommand{\ownNineChannelRS}{2.0\%}
\newcommand{\deltaNineChannelRSmFLN}{0.5}
\newcommand{\mcMedianOwnership}{0.0\%}
\newcommand{\mcMeanOwnership}{0.0\%}
\newcommand{\mcLowCIOwnership}{0.0\%}
\newcommand{\mcHighCIOwnership}{0.0\%}
\newcommand{\mcLowIQROwnership}{0.0\%}
\newcommand{\mcHighIQROwnership}{0.0\%}
\newcommand{\mcMinOwnership}{0.0\%}
\newcommand{\mcMaxOwnership}{0.0\%}
\newcommand{\nMCDraws}{1,000}
\newcommand{\pMwrLoadNum}{13}
\newcommand{\pctHRSAboveWmaxNum}{0.4}
\newcommand{\pctHRSLifetimeNum}{3.11}
\newcommand{\pctHRSIannPooledNum}{5.21}
\newcommand{\pctHRSLifetimeCILowNum}{2.54}
\newcommand{\pctHRSLifetimeCIHighNum}{3.81}
\newcommand{\pctHRSIannCILowNum}{4.45}
\newcommand{\pctHRSIannCIHighNum}{6.09}
\newcommand{\pctELSAWaveSixAnnuityNum}{90.2}
\newcommand{\pctELSAPostAnnuityNum}{1.27}
\newcommand{\ownAddSSNum}{100.0}
\newcommand{\ownAddBequestsNum}{100.0}
\newcommand{\ownAddMedRSNum}{87.7}
\newcommand{\ownAddPessimismNum}{62.4}
\newcommand{\ownAddLoadsNum}{18.2}
\newcommand{\ownSixChannelNum}{8.5}
\newcommand{\ownSevenChannelExtNum}{2.6}
\newcommand{\ownEightChannelExtNum}{1.8}
\newcommand{\ownNineChannelLTCNum}{2.0}
\newcommand{\ownNoBehavioralBaselineNum}{2.0}
\newcommand{\ownTenChannelSDUNum}{51.1}
\newcommand{\ownElevenChannelFullNum}{0.0}
\newcommand{\ownModelOneNum}{0.0}
\newcommand{\ownModelOneStructuralNum}{0.0}
\newcommand{\pUKRetentionPreNum}{95}
\newcommand{\pUKRetentionLowNum}{13}
\newcommand{\pUKRetentionMidNum}{17}
\newcommand{\pUKRetentionHighNum}{25}
\newcommand{\ownWedgeLowNum}{5.7}
\newcommand{\ownWedgeMidNum}{7.5}
\newcommand{\ownWedgeHighNum}{11.0}
\newcommand{\ownModelTwoLowNum}{5.7}
\newcommand{\ownModelTwoMidNum}{7.5}
\newcommand{\ownModelTwoHighNum}{11.0}
\newcommand{\ownModelTwoNum}{7.5}
\newcommand{\ownHeadlineNum}{7.5}
\newcommand{\ownBracketLowNum}{5.7}
\newcommand{\ownBracketHighNum}{11.0}
\newcommand{\retentionSSNum}{239.0}
\newcommand{\ownPrePessimismNum}{87.7}
\newcommand{\retentionPessimismNum}{71.1}
\newcommand{\retentionLoadsNum}{29.1}
\newcommand{\retentionInflationNum}{46.6}
\newcommand{\retentionMedRSNum}{87.7}
\newcommand{\shapShareSSNum}{-27}
\newcommand{\shapShareBequestsNum}{1}
\newcommand{\shapShareMedRSNum}{37}
\newcommand{\shapSharePessimismNum}{13}
\newcommand{\shapShareAgeNeedsNum}{6}
\newcommand{\shapShareStateUtilNum}{1}
\newcommand{\shapShareLoadsNum}{39}
\newcommand{\shapShareInflationNum}{3}
\newcommand{\shapShareLTCNum}{1}
\newcommand{\shapShareSDUNum}{-73}
\newcommand{\shapSharePEDNum}{100}
\newcommand{\shapShareRSPessAgeNum}{56}
\newcommand{\shapShareNineSSNum}{-62}
\newcommand{\shapShareNineBequestsNum}{14}
\newcommand{\shapShareNineMedRSNum}{29}
\newcommand{\shapShareNinePessimismNum}{29}
\newcommand{\shapShareNineAgeNeedsNum}{12}
\newcommand{\shapShareNineStateUtilNum}{1}
\newcommand{\shapShareNineLoadsNum}{67}
\newcommand{\shapShareNineInflationNum}{11}
\newcommand{\shapShareNineLTCNum}{0}
\newcommand{\cevGoodHundredKNoBqNum}{0.0}
\newcommand{\cevGoodHundredKDFJNum}{0.0}
\newcommand{\cevGoodHundredKStrongNum}{0.0}
\newcommand{\cevGoodTwoHundredKNoBqNum}{0.0}
\newcommand{\cevGoodTwoHundredKDFJNum}{0.0}
\newcommand{\cevGoodTwoHundredKStrongNum}{0.0}
\newcommand{\cevGoodOneMillNoBqNum}{6.6}
\newcommand{\cevGoodOneMillDFJNum}{0.8}
\newcommand{\cevGoodOneMillStrongNum}{0.0}
\newcommand{\popMeanCevNoBqNum}{0.5}
\newcommand{\popMeanCevDFJNum}{0.0}
\newcommand{\popFracPosNoBqNum}{15.8}
\newcommand{\popFracPosDFJNum}{7.2}
\newcommand{\popFracAboveOneNoBqNum}{11.8}
\newcommand{\popFracAboveOneDFJNum}{0.0}
\newcommand{\ownGroupPricingNum}{6.1}
\newcommand{\ownPublicOptionNum}{14.3}
\newcommand{\ownActuariallyFairNum}{25.6}
\newcommand{\ownRealTIPSNum}{0.0}
\newcommand{\ownRealNomEquivNum}{8.1}
\newcommand{\ownFairRealNum}{31.8}
\newcommand{\ownCorrectPessimismNum}{18.3}
\newcommand{\ownGroupPlusCorrectNum}{25.9}
\newcommand{\ownBestFeasibleNum}{35.0}
\newcommand{\cevGoodHundredKBaselineNum}{0.0}
\newcommand{\cevGoodHundredKGroupPriceNum}{0.0}
\newcommand{\cevGoodHundredKRealAnnNum}{0.0}
\newcommand{\cevGoodHundredKBestFeasNum}{0.0}
\newcommand{\cevGoodTwoHundredKBaselineNum}{0.0}
\newcommand{\cevGoodTwoHundredKGroupPriceNum}{0.0}
\newcommand{\cevGoodTwoHundredKRealAnnNum}{0.0}
\newcommand{\cevGoodTwoHundredKBestFeasNum}{0.7}
\newcommand{\cevGoodFiveHundredKBaselineNum}{0.0}
\newcommand{\cevGoodFiveHundredKGroupPriceNum}{0.5}
\newcommand{\cevGoodFiveHundredKRealAnnNum}{0.1}
\newcommand{\cevGoodFiveHundredKBestFeasNum}{4.1}
\newcommand{\cevGoodOneMillBaselineNum}{0.8}
\newcommand{\cevGoodOneMillGroupPriceNum}{1.5}
\newcommand{\cevGoodOneMillRealAnnNum}{1.0}
\newcommand{\cevGoodOneMillBestFeasNum}{5.6}
\newcommand{\ownGammaTwoNum}{0.0}
\newcommand{\ownGammaTwoPointThreeNum}{0.0}
\newcommand{\ownGammaTwoPointFourNum}{0.1}
\newcommand{\ownGammaTwoPointFiveNum}{1.8}
\newcommand{\ownGammaThreeNum}{18.1}
\newcommand{\ownInflationOneNum}{4.7}
\newcommand{\ownInflationTwoNum}{1.8}
\newcommand{\ownInflationThreeNum}{0.7}
\newcommand{\ownPsiObjectiveNum}{21.4}
\newcommand{\ownPsiNinetySevenNum}{5.1}
\newcommand{\ownPsiBaselineNum}{8.7}
\newcommand{\ownPsiNinetyNineNum}{14.0}
\newcommand{\ownMWREightyTwoNum}{0.0}
\newcommand{\ownMWREightyFiveNum}{0.4}
\newcommand{\ownMWRNinetyNum}{5.8}
\newcommand{\ownMWRNinetyFiveNum}{15.4}
\newcommand{\ownBequestNoneNum}{21.0}
\newcommand{\ownHazardRSNum}{3.2}
\newcommand{\ownHazardHRSNum}{6.4}
\newcommand{\ownHazardConservativeNum}{10.6}
\newcommand{\ownHazardAgeBandNum}{7.3}
\newcommand{\ownSSCutZeroNum}{1.8}
\newcommand{\ownSSCutTenNum}{3.8}
\newcommand{\ownSSCutFifteenNum}{5.9}
\newcommand{\ownSSCutTwentyThreeNum}{10.4}
\newcommand{\ownSSCutThirtyNum}{16.3}
\newcommand{\ownSSCutFortyNum}{25.1}
\newcommand{\ownSSCutFiftyNum}{28.6}
\newcommand{\ownSSCutHundredNum}{14.3}
\newcommand{\ownNineChannelFLNNum}{1.5}
\newcommand{\ownNineChannelRSNum}{2.0}
\newcommand{\mcMedianOwnershipNum}{0.0}
\newcommand{\mcMeanOwnershipNum}{0.0}
\newcommand{\mcLowCIOwnershipNum}{0.0}
\newcommand{\mcHighCIOwnershipNum}{0.0}
\newcommand{\mcLowIQROwnershipNum}{0.0}
\newcommand{\mcHighIQROwnershipNum}{0.0}
\newcommand{\mcMinOwnershipNum}{0.0}
\newcommand{\mcMaxOwnershipNum}{0.0}


% Defensive placeholders. \providecommand only defines a macro if it isn't
% already defined in numbers.tex, so the appendix compiles cleanly during
% the transition period before the pipeline regenerates numbers.tex.
\providecommand{\ownNoBehavioralBaseline}{[pending]}
\providecommand{\ownTenChannelSDU}{[pending]}
\providecommand{\ownNineChannelFull}{\ownNoBehavioralBaseline}
\providecommand{\ownNineChannelFullNum}{\ownNoBehavioralBaselineNum}
\providecommand{\ownNineChannelSDU}{\ownTenChannelSDU}
\providecommand{\pLambdaW}{0.625}
\providecommand{\pPsiPurchase}{0.05}
\providecommand{\pChiLTC}{0.49}

\title{Online Appendix for ``Loads, Not Bequests: An Order-Independent Decomposition of the Annuity Puzzle''}
\author{Derek Tharp}
\date{\today}

\begin{document}

\maketitle

\appendix
\setcounter{section}{0}
\renewcommand{\thesection}{\Alph{section}}
\renewcommand{\thetable}{A\arabic{table}}
\renewcommand{\thefigure}{A\arabic{figure}}


%% ================================================================
%%  A. MODEL DETAILS
%% ================================================================
\section{Model Details}
\label{app:model}

\subsection{Health transition matrices}
\label{app:health_transitions}

Health transitions are calibrated to the RAND HRS longitudinal file (consecutive-wave pairs from waves 5--12, 2000--2014, with two-year rates annualized via matrix square root), with the five self-reported health categories mapped to three states: Good (Excellent/Very Good), Fair (Good), and Poor (Fair/Poor). Transition probabilities are age-band dependent: distinct $3 \times 3$ transition matrices are estimated for six age bands (65--69, 70--74, 75--79, 80--84, 85--89, 90+) and applied piecewise within each band. The age-band specification captures the empirical acceleration in poor-health entry around ages 75--85 (Hurd-Michaud-Rohwedder 2017) more accurately than a two-anchor linear interpolation. Each band's transition matrix is row-normalized after estimation to ensure exact stochasticity. Reference values follow \citet{denardi2010}.

The transition matrices exhibit two standard features. First, health deterioration is more likely than improvement: transition rates from Good to Fair exceed those from Fair to Good in every age band from 70--74 onward (in the 65--69 band the two rates are nearly equal, 0.144 versus 0.147). Second, deterioration accelerates with age: the annual probability of remaining in Good health falls from 0.84 at ages 65--69 to 0.67 at ages 90+.

\subsection{Medical expenditure calibration}

Log out-of-pocket medical expenditures follow:
\[
\ln m_t = \mu_m^{\text{base}} + g \cdot (t - 65) + \delta_\mu(H_t) + \big[\sigma_m^{\text{base}} + \delta_\sigma(H_t)\big] \cdot \varepsilon_t, \quad \varepsilon_t \sim N(0,1),
\]
where $\mu_m^{\text{base}} = 7.037$ is the log mean at age 65 for Fair health, $g = 0.0652$ is the annual growth rate in the log mean, $\delta_\mu(H)$ shifts the log mean by health state $[-0.5, 0, 0.7]$ for Good, Fair, and Poor respectively, $\sigma_m^{\text{base}} = 1.4$ is the base log standard deviation, and $\delta_\sigma(H)$ shifts the log standard deviation by health $[-0.2, 0, 0.2]$.

\begin{table}[htbp]
\centering
\caption{Medical Expenditure Moments: Model vs.\ Target}
\label{tab:medical_moments}
\begin{tabular}{lrr}
\toprule
Moment & Model & Target (Jones et al.\ 2018) \\
\midrule
Mean OOP, age 70, Fair health & \medMeanSeventy & \$4,200 \\
Mean OOP+Medicaid, age 100, Fair health & \medMeanHundred & \$29,700 \\
95th pctile OOP+Medicaid, age 100, Fair health & \medPctNinetyFive & \$111,200 \\
\bottomrule
\end{tabular}
\end{table}

The lognormal specification produces a right-skewed distribution consistent with the heavy tails observed in medical expenditure data. The log standard deviation of 1.4 implies a coefficient of variation of approximately 2.5, which matches the empirical distribution of out-of-pocket spending in the HRS and MEPS.

\subsection{Survival calibration}
\label{app:survival}

Baseline survival probabilities $s_{\text{base}}(t)$ use a sex-blended SSA 2003 period life table: the survivorship mixture of the male and female SSA 2003 tables at the model-eligible sample's female share (\pctFemaleShare), $S_{\text{base}}(a) = w_f S_f(a) + (1-w_f) S_m(a)$, with one-year conditional survivals derived from the mixed survivorship. The prior convention---the male table from the \citet{lockwood2012} replication code (BAP\_wtp.m)---is retained as the mortality-convention robustness of Appendix~\ref{app:male_mortality}.

Health-conditional survival is anchored to this aggregate table. The raw proportional-hazard form is $s_{\text{base}}(t)^{\mu(H)}$, but the multipliers are relative hazards rather than table-anchored levels, so their population-weighted average departs from $s_{\text{base}}(t)$. We rescale every multiplier by a common age-specific factor $a_t$ chosen so the health mixture reproduces the table exactly:
\[
s(t, H) = s_{\text{base}}(t)^{a_t\,\mu(H)}, \qquad a_t \text{ solves } \sum_H \pi_H(t)\, s_{\text{base}}(t)^{a_t\,\mu(H)} = s_{\text{base}}(t),
\]
where $\pi_H(t)$ is the health distribution among survivors at age $t$, initialized from the sample's pooled ages 65--69 shares (957 Good, 761 Fair, 561 Poor of 2{,}279) and evolved forward under the health transition matrices weighted by the normalized survivals. The factor $a_t$ is unique (the left side is strictly decreasing in $a_t$) and solved by bisection. Hazard ratios across health states are preserved; only the level is anchored. The multipliers $[\pHazardGood, \pHazardFair, \pHazardPoor]$ for Good, Fair, and Poor health are informed by age-band estimates from the RAND HRS longitudinal file:

\begin{center}
\begin{tabular}{lccc}
\toprule
Age band & $\mu(\text{Good})$ & $\mu(\text{Fair})$ & $\mu(\text{Poor})$ \\
\midrule
65--74 (midpoint 69.5) & 0.49 & 1.00 & 3.29 \\
75--84 (midpoint 79.5) & 0.60 & 1.00 & 2.77 \\
85+ (midpoint 90.0) & 0.74 & 1.00 & 1.82 \\
\bottomrule
\end{tabular}
\end{center}

The constant baseline multipliers $[\pHazardGood, \pHazardFair, \pHazardPoor]$ sit at the upper end of these age-band estimates: the Good and Fair multipliers fall between the HRS self-reported values and the R-S functional estimates, while the Poor multiplier sits just above both, reflecting the steeper functional-limitation gradient at the younger ages where the annuitization decision is made. Robustness checks in Section~\ref{M-sec:robustness} and Appendix~\ref{app:robustness} span the full range from $[0.45, 1.0, 3.5]$ to $[0.60, 1.0, 2.0]$. The code also supports age-varying multipliers with linear interpolation between band midpoints, available as a further extension.

\subsection{Subjective survival beliefs}
\label{app:survival_pessimism}

The agent uses subjective survival probabilities in the Bellman equation:
\[
s^{\text{subj}}(t, H) = \psi \cdot s(t, H), \quad \psi \in (0, 1].
\]
Annuity pricing uses objective population survival (the population life table, not health-rated), creating a wedge: the agent perceives annuities as less valuable than they are because she overweights the probability of dying before collecting.

\citet{odeasturrock2023} estimate that individuals aged 65--69 underestimate survival to 75 by about 15 percentage points; their ten-year ratio gives the per-year factor $\psi = (0.71/0.86)^{1/10} \approx \pPessimism$, the baseline value. The baseline is the paper's focal empirical translation of the age-matched evidence, not an endpoint chosen for fit. The direction of survival miscalibration is age-dependent: \citet{heimer2019} document a crossover near ages 65--70, after which individuals tend to \emph{overestimate} survival, so for this cohort mild optimism ($\psi$ at or above one) is also plausible. The paper therefore reports $\psi$ over the range $[0.96, 1.0]$---a strong-pessimism endpoint (at which subjective 10-year survival is $0.96^{10} \approx 0.66$ times objective), the O'Dea--Sturrock baseline, and objective beliefs. The participation level is sensitive to this choice (Section~\ref{M-sec:robustness}); the channel ranking is not (Appendix~\ref{app:psi_ranking}).

\subsection{Channel ranking at the strong-pessimism endpoint}
\label{app:psi_ranking}

Because the participation level is sensitive to $\psi$, a natural question is whether the channel ranking is too. It is not. Recomputing the full nine-channel Shapley game at the strong-pessimism endpoint $\psi = 0.96$ (Table~\ref{tab:shapley_psi_endpoint}) leaves pricing loads the dominant demand suppressor and bequest motives mid-pack---the two ranking claims the paper relies on. Survival pessimism's own Shapley contribution is larger at the endpoint than at the baseline, as expected when beliefs are made more pessimistic (\psiEndpointPessimismAbs{} versus \psiEndpointNextAbs~pp for the next-largest suppressor), and it holds the second rank at both ends of the range. The decomposition is computed on the $60 \times 20 \times 101$ grid, which the convergence table places within about two percentage points of production (Appendix~\ref{app:grid_convergence}); the level is not the object here.

\input{../tables/tex/shapley_psi_endpoint.tex}

\subsection{Age-varying consumption needs}
\label{app:age_needs}

The age-varying needs weight $w(t) = (1 - \delta_c)^{t-1}$ with $\delta_c = \pDeltaC$ captures the declining consumption expenditure profile documented by \citet{aguiarhurst2013}. At this calibration, the weight declines from 1.0 at age 65 to 0.67 at age 85 and 0.49 at age 100. This implies that the felicity from a given consumption level at age 85 is roughly two-thirds of the felicity at age 65.

The channel reduces annuity demand because the annuity delivers a level income stream while the individual's consumption needs decline. An agent with declining needs prefers front-loaded consumption, which liquid wealth supports better than an annuity stream. The Shapley value of \shapNineAgeNeeds{} pp (\shapShareNineAgeNeeds{} share) indicates this is a quantitatively meaningful channel, comparable to the survival-pessimism contribution (\shapNinePessimism{} pp) and far larger than the near-zero inflation-erosion contribution.


\subsection{State-dependent utility calibration}
\label{app:state_dependent}

The health-utility weights $\varphi = [\pHealthUtilGood, \pHealthUtilFair, \pHealthUtilPoor]$ are based on the central estimates of \citet{finkelsteinluttmer2013}, who used variation in health induced by the onset of chronic conditions to estimate the effect of health on the marginal utility of consumption. Their central estimates imply that individuals in poor health value a dollar of consumption roughly 10--25\% less than those in good health. \citet{reichlingsmetters2015} adopted a softer translation of these estimates, $\varphi = [1.00, \pHealthUtilFairRS, \pHealthUtilPoorRS]$, noting that HRS self-reported health categories do not map perfectly onto FLN's chronic-condition onset states and that FLN's confidence intervals cover the full 10--25\% range.

The production calibration uses the FLN central midpoint $\varphi = [\pHealthUtilGood, \pHealthUtilFair, \pHealthUtilPoor]$, under which the full eight-channel rational+preference model predicts \ownEightChannelExt{} ownership. The state-utility mapping moves this prediction only modestly: the harsher FLN raw endpoint yields \ownEightChFLNRaw{} and the softer Reichling--Smetters mapping \ownNineChannelRS{}, a difference of \deltaNineChannelRSmFLN{} pp. The Reichling--Smetters mapping is reported as a robustness check. In the calibrated model, state-dependent utility has a negligible effect on predicted ownership (Shapley value of \shapNineStateUtil{} pp) in either calibration. This near-zero effect reflects the modest magnitude of the health-utility weights and the fact that poor-health states already receive low weight through the mortality channel: individuals in Poor health face high mortality, so their contribution to the annuity's expected value is small regardless of the utility weight.

\subsection{Behavioral parameter sensitivity}
\label{app:behavioral_sensitivity}

The behavioral parameters $\lambda_W$ and $\psi_{\text{purchase}}$ are not separately identified from a single bundling moment and are not moment-matched within the structural model. They are reported as \emph{exploratory} parameters anchored on literature magnitudes. The point values used for illustrative within-model robustness exercises are $\lambda_W = \pLambdaW$ (the Blanchett--Finke point estimate of the source-dependent utility spending differential between guaranteed income and portfolio wealth; \citealp{blanchett2024,blanchett2025}) and $\psi_{\text{purchase}} = \pPsiPurchase$ (a literature-magnitude best guess for narrow-framing purchase-event disutility consistent with the order of magnitudes documented in \citealp{barberishuang2009} and \citealp{brown2008framing}).

Within-model sensitivity is reported over the ranges $\lambda_W \in \{0.5, 0.625, 0.75, 0.85\}$ and $\psi_{\text{purchase}} \in \{0.01, 0.05, 0.09\}$. These ranges are stated for transparency about the exploratory status of the parameters: any single combination within the grid is consistent with the available behavioral evidence at the literature's order of magnitude. The body of the paper does not rely on a specific point value of either parameter for its headline result. The disciplined reading is the nine-channel structural baseline of \ownNineChannelLTC{}.

The HRS empirical targets reported in the body for out-of-sample comparison are $\pctHRSLifetime$ (95\% CI $[\pctHRSLifetimeCILow, \pctHRSLifetimeCIHigh]$, lifetime contract indicator) and $\pctHRSIannPooled$ (95\% CI $[\pctHRSIannCILow, \pctHRSIannCIHigh]$, conventional any-annuity income proxy).

\paragraph{Extension path.} Table~\ref{tab:extension_path} tabulates the sequential extension path from the six rational channels through the nine-channel structural baseline to the exploratory eleven-channel specification, isolating each preference-channel increment and the two behavioral increments discussed in the main text.

\begin{table}[htbp]
\centering
\caption{Two-Model Decomposition: Structural Channels and UK Reduced-Form Transport}
\label{tab:extension_path}
\begin{threeparttable}
\begin{tabular}{lcc}
\toprule
Specification & Ownership (\%) & $\Delta$ (pp) \\
\midrule
Six rational channels (Layer~1)              & 8.5 & --- \\
+ Age-varying consumption needs              & 2.6 & -5.8 \\
+ State-dependent utility                    & 1.8 & -0.9 \\
+ Public-care aversion $\chi_{\text{LTC}}$ (Layer~2 complete)         & 1.7 & -0.1 \\
+ Source-dependent utility (Force A)         & 27.8 & +26.1 \\
+ Narrow-framing penalty (Force B; Model 1) & 0.0 & -27.8 \\
\midrule
Model 2: frictionless baseline (41.8) $\times$ UK 17/95 & 7.5 & --- \\
\bottomrule
\end{tabular}
\begin{tablenotes}
\small
\item Model 1 (top block) is the structural multi-channel decomposition.
Layer 1 covers rational frictions; Layer 2 adds preference and structural
channels; Force A (SDU) and Force B (PED) close the model with behavioral
channels calibrated to Blanchett-Finke (2024-25) and Chalmers-Reuter (2012)
respectively.
\item Model 2 (bottom row) applies the UK 2015 pension-freedoms
\item retention factor (UK post / UK pre = 0.18 at the production midpoint) to the
\item frictionless Yaari baseline. The UK retention range [13\%, 25\%] maps to a Model~2 prediction bracket of [5.7\%, 11.0\%].
\end{tablenotes}
\end{threeparttable}
\end{table}


\paragraph{Eleven-channel extended Shapley.} Table~\ref{tab:shapley_exact} reports the exact Shapley decomposition overall $2^{11} = 2{,}048$ subsets of the eleven-channel game, which adds the two behavioral parameters to the nine structural channels. Under the literature-magnitude parameterization the narrow-framing penalty carries the largest absolute contribution (\shapPED{} pp, \shapSharePED) and source-dependent utility a large opposite-sign booster contribution (\shapSDU{} pp, \shapShareSDU). These magnitudes are artifacts of the chosen, un-identified parameter values rather than measured forces: at $\psi_{\text{purchase}} = \pPsiPurchase$ the at-purchase penalty saturates the model, eliminating participation for nearly all agents, so a near-total Shapley share is mechanical, and the one-decimal values for the two behavioral rows should be read as sign-definite and large rather than as interpretable point magnitudes. The behavioral players also rescale the structural attributions relative to the nine-channel game rather than shrinking them uniformly: pricing loads fall from \shapNineLoads{} pp in the structural game to \shapLoads{} pp here yet remain the top structural suppressor, survival pessimism falls to \shapPessimism{} pp, and the Med+R-S channel rises to \shapMedRS{} pp, moving up the lower tiers. The behavioral channels dilute the structural attribution without reversing its top; the lower-tier shifts are a saturation artifact---with the grand coalition near zero ownership, coalition margins are dominated by near-zero-participation subsets, which favor the channels that bind there---and that artifact is precisely why the structural ranking is read off the nine-channel game and this table is presented only as an illustration of potential behavioral scale, not as an identified attribution.

\begin{table}[htbp]
\centering
\caption{Exact Shapley-Value Decomposition of Predicted Annuity Ownership}
\label{tab:shapley_exact}
\begin{tabular}{lcc}
\toprule
Channel & Shapley (pp) & Share (\%) \\
\midrule
SS & -17.65 & -42.2 \\
Bequests & +0.85 & 2.0 \\
Medical+R-S & +8.60 & 20.6 \\
Pessimism & +5.85 & 14.0 \\
Age needs & +2.90 & 6.9 \\
State utility & +0.26 & 0.6 \\
Loads & +14.01 & 33.5 \\
Inflation & +1.44 & 3.4 \\
Public-care aversion (LTC) & -0.64 & -1.5 \\
Source-dependent utility (SDU) & -13.44 & -32.1 \\
Narrow-framing penalty (PED) & +39.65 & 94.7 \\
\midrule
Total & +41.85 & 100.0 \\
\bottomrule
\end{tabular}
\begin{tablenotes}
\small
\item Exact Shapley values computed from all 2048 channel subsets.
Each value represents the weighted average marginal ownership reduction (pp)
across all coalition orderings.
Yaari baseline: 41.8\%. Full model: 0.0\%.
\end{tablenotes}
\end{table}


\subsection{Nominal annuity pricing derivation}
\label{app:nominal_pricing}

An insurer selling a nominal annuity invests the premium in nominal bonds yielding approximately $r_{\text{nom}} = (1+r)(1+\pi) - 1$ (exact Fisher equation). The present value of the annuity payment stream, discounted at the nominal rate, is:
\[
\text{PV} = \sum_{k=0}^{T-1} \frac{\prod_{j=0}^{k-1} s_{\text{base}}(j)}{(1+r_{\text{nom}})^k}.
\]
The fair payout rate is $1/\text{PV}$. The loaded payout rate is $\text{MWR}/\text{PV}$.

The consumer receives nominal payments $A_1$ per period. In real terms, the payment in period $t$ is $A_1 / (1+\pi)^{t-1}$. The initial real payout under nominal pricing exceeds the payout under real pricing (which uses discount rate $r$), but declines in purchasing power over time.

When $\pi = 0$, $r_{\text{nom}} = r$ and the nominal and real pricing formulas coincide. This ensures that the Yaari benchmark (which assumes $\pi = 0$) is unaffected by the pricing convention.


%% ================================================================
%%  B. GRID CONVERGENCE
%% ================================================================
\section{Grid Convergence}
\label{app:grid_convergence}

Table~\ref{tab:grid_convergence} reports predicted ownership and mean annuitization share ($\bar{\alpha}$) under the full nine-channel model, computed on the headline per-quartile Social Security configuration so that the production cell reproduces the headline ownership level rather than a mean-income proxy. Panel A holds the quadrature rule at 9-node Gauss--Hermite and varies the wealth--annuity grid; Panel B holds the grid at $80 \times 30$ and varies the quadrature node count; Panel C reports the most-refined reference; Panel D varies the annuitization grid $n_\alpha$.

\input{../tables/tex/grid_convergence.tex}

Because the diagnostic uses the production per-quartile assignment, the production cell ($80 \times 30$, 9-node) reproduces the headline ownership of \ownNineChannelLTC{} exactly. Holding quadrature at the production 9 nodes (Panel A), grid refinement is non-monotone: ownership is \gridConvMed{} at $60 \times 20$, \gridConvProd{} at the production $80 \times 30$, \gridConvFine{} at $100 \times 40$, and \gridConvVeryFine{} at $120 \times 50$ (the change between the two finest grids is \gridConvFineStep{}~pp), so the finest grid sits about \gridConvBias{}~pp below production. This non-monotonicity reflects the discrete extensive margin near the fixed-cost threshold---the same level fragility documented in Section~\ref{M-sec:robustness}---interacting with quadrature error rather than incomplete grid resolution: raising the finest grid to the 15-node reference (Panel C) returns ownership to \gridConvRef{}, within \gridConvTotalUnc{}~pp of the headline.

Quadrature refinement is tighter. At the production $80 \times 30$ grid (Panel B), ownership at $n_{\text{quad}} \in \{9, 11, 13, 15\}$ is \quadNine{}, \quadEleven{}, \quadThirteen{}, and \quadFifteen{}---a \quadBand{}~pp band, stable rather than monotone. Lower-order rules ($n_{\text{quad}} \leq 7$) deviate more, by up to \quadLowDeviation{}~pp, because the lognormal medical expense distribution ($\sigma \approx 1.4$) has heavy right tails that interact with the Medicaid floor; insufficient tail resolution at low orders shifts marginal households across the discrete ownership threshold ($\alpha^* > 0$).

Nine-node Gauss--Hermite is adopted as the production standard, balancing accuracy against computational cost; taking the most-refined specification (Panel C: $120 \times 50$ grid, 15-node GH) as reference, the total numerical uncertainty---combining grid coarseness and quadrature error---is approximately \gridConvTotalUnc{}~pp. Panel D varies the age-65 annuitization grid $n_\alpha$ over $\{51, 101, 201, 401\}$ while holding the wealth grid and quadrature at production resolution; because ownership is the discontinuous indicator $\alpha^* > 0$, this checks that the participation level is not pinned to the grid's coarseness. Predicted ownership stays within a \alphaGridBand{}~pp band (\alphaGridCoarse{} to \alphaGridVeryFine{}, \alphaGridProd{} at the production $n_\alpha = \pNAlpha$). The top of the channel ranking is stable across all grid and quadrature specifications tested, which is the property the paper's conclusions rely on. The numerical variation is small relative to the leading Shapley contributions (pricing loads and survival pessimism) but not relative to the minor channels (Med+R-S, state-dependent utility, inflation, public-care aversion), whose small attributions are not sharply resolved at this grid resolution.


%% ================================================================
%%  C. BEQUEST SPECIFICATION DETAILS
%% ================================================================
\section{Bequest Specification}
\label{app:bequest}

The De Nardi--French--Jones (DFJ) bequest utility is $v(b) = \theta (b + \kappa)^{1-\gamma}/(1-\gamma)$. The parameter $\kappa$ controls the luxury-good nature of bequests. When $\kappa$ is large relative to wealth $b$, marginal bequest utility $v'(b) = \theta(b + \kappa)^{-\gamma}$ is nearly constant across wealth levels, making bequests a luxury good: the bequest-to-wealth ratio rises with wealth.

A numerical floor of \$1 is applied to $b + \kappa$ to avoid undefined values when $\kappa = 0$ and $b = 0$. This floor does not bind under the DFJ specification ($\kappa = \pKappaDFJ$) but affects homothetic counterfactuals near zero bequests.

The reference parameters ($\theta = \pThetaDFJ$, $\kappa = \pKappaDFJ$) are from \citet{lockwood2012}, who estimated them from HRS data under the DFJ specification at $\sigma = 2$ (his notation for CRRA). These values are used at all $\gamma$ in the present model without recalibration (see Section~\ref{app:bequest_recal} for discussion). At these values, an individual with \$50,000 in wealth has $v'(b) \approx \theta \cdot (322{,}628)^{-\gamma}$, which is negligible compared to the marginal utility of consumption. The bequest motive binds meaningfully only for individuals with wealth approaching $\kappa$.

\paragraph{Homothetic bequest anomaly.} Combining the DFJ bequest intensity $\theta = \pThetaDFJ$ with a homothetic specification ($\kappa = \$10$) produces anomalous results. At $\kappa = \$10$, marginal bequest utility near $b = 0$ is $\theta \cdot 10^{-2.5} = 56.96 \times 3.16 \times 10^{-3} \approx 0.18$, which is extreme relative to marginal consumption utility. This creates a narrow precautionary buffer---individuals avoid annuitizing not to protect bequests, but to avoid the near-zero-wealth states where $v'(b)$ diverges. The predicted ownership under this misspecification is higher than the no-bequest case, because the buffer effect dominates the bequest effect.

This anomaly illustrates why $\theta$ and $\kappa$ must be interpreted as a pair, not independently. The DFJ estimates from \citet{lockwood2012} are jointly calibrated; using one with the other's value replaced is not a valid counterfactual.

\subsection{Bequest parameter portability}
\label{app:bequest_recal}

\citet{lockwood2012} estimated $\theta = \pThetaDFJ$ and $\kappa = \pKappaDFJ$ at $\sigma = 2$ (his notation for the CRRA coefficient). When CRRA changes from $\gamma_{\text{ref}}$ to $\gamma_{\text{new}}$, one could recalibrate $\theta$ to maintain the same optimal bequest-to-consumption ratio via
\[
\theta_{\text{new}} = \theta_{\text{ref}}^{\gamma_{\text{new}}/\gamma_{\text{ref}}}.
\]
At $\gamma = \pGamma$, this would yield $\theta = \pThetaDFJ^{\pGamma/2.0} = 156.5$, nearly tripling the bequest weight. However, this recalibration is ad hoc: Lockwood's $\theta$ was estimated from data, not derived from first principles, and there is no guarantee that the first-order condition used to derive the scaling formula holds in the full model with medical expenditure risk, stochastic mortality, and survival pessimism.

Table~\ref{tab:bequest_recal} reports a forward simulation test at $\gamma = 2.0$ and $\gamma = \pGamma$ using the original $\theta = \pThetaDFJ$. Holding $\theta$ fixed while raising $\gamma$ from Lockwood's estimation value to the baseline leaves the realized bequest-to-wealth ratio modestly \emph{above} his estimation point (Table~\ref{tab:bequest_recal}, divergence below the 20\% retargeting threshold), so the production calibration already affords the bequest motive slightly more weight than a retargeting to his moment would, not less. The analytic portability formula pushes the other way, raising $\theta$ to $156.5$ and strengthening the marginal bequest incentive further; the two adjustments bracket the production value and disagree in direction, which is why neither is imposed and all results use Lockwood's original $\theta$ at every $\gamma$. The conclusion the paper draws from the bequest channel concerns its \emph{rank}, not its level, and the rank is stable across the entire risk-aversion sweep at fixed $\theta$: pricing loads are the leading suppressor on the ownership statistic at every $\gamma \in [2.0, 3.0]$, and bequest motives, although they rise from mid-pack at the baseline into the top three as risk aversion falls, overtake loads only at $\gamma = 1.5$, the degenerate zero-ownership boundary (Section~\ref{M-sec:robustness}). Because that sweep already traverses a wide range of effective bequest strength without displacing pricing loads, a recalibration of $\theta$ within the empirically relevant range would not overturn the ranking. The predicted ownership \emph{level} is sensitive to the bequest calibration; the channel ranking is not.

% Generated by calibration/recalibrate_bequests.jl
\begin{table}[htbp]
\centering
\caption{Bequest Parameter Portability: $\gamma = 2.0$ vs $\gamma = 2.5$}
\label{tab:bequest_recal}
\begin{tabular}{lcc}
\toprule
 & $\gamma = 2.0$ & $\gamma = 2.5$ \\
\midrule
$\theta$ (bequest intensity) & 56.96 & 56.96 \\
$\kappa$ (bequest shifter) & \$272,628 & \$272,628 \\
Bequest / initial wealth & 0.163 & 0.185 \\
Mean bequest & \$40765 & \$46336 \\
Divergence & \multicolumn{2}{c}{13.7\%} \\
\bottomrule
\end{tabular}
\begin{tablenotes}
\small
\item Both columns use original $\theta = 56.96$ from Lockwood (2012) at $\sigma = 2$.
Divergence in bequest-to-wealth ratio is 13.7\%, below the threshold that would
justify the strong assumptions embedded in analytical recalibration (see text).
Simulated 50000 trajectories, initial wealth \$250000, Fair health.
\end{tablenotes}
\end{table}



%% ================================================================
%%  D. PASHCHENKO COMPARISON DETAILS
%% ================================================================
\section{Comparison with Pashchenko (2013)}
\label{app:pashchenko}

\citet{pashchenko2013} conducted a sequential decomposition with four channels: pre-existing annuitization (Social Security and DB pensions), bequest motives, minimum purchase requirements, and pricing loads. Her full model predicted approximately 20\% participation, roughly four times the observed rate of 5\%.

Table~\ref{tab:pashchenko} evaluates her channel set within the present framework. This comparison differs from a replication in three respects: the present model uses DFJ luxury-good bequests rather than the homothetic specification she employed; a continuous annuitization fraction rather than a binary purchase decision; and no housing illiquidity. The six-channel rational model under the present reformulation predicts \ownSixChannel{}. The two are not directly comparable in level---her model omits the combined Med+R-S correlation, survival pessimism, age-varying needs, and inflation erosion; the present model omits her housing illiquidity and minimum purchase channels---but the qualitative finding is preserved: standard rational channels alone substantially overshoot observed US ownership. Adding the two preference channels (age-varying consumption needs and state-dependent utility) reduces the present prediction to \ownEightChannelExt{}, and adding the structural public-care aversion channel yields the nine-channel structural baseline of \ownNineChannelLTC{}, which is the disciplined reading.

\begin{table}[htbp]
\centering
\caption{Pashchenko (2013) Channels in Our Framework}
\label{tab:pashchenko}
\begin{tabular}{lcc}
\toprule
Model Specification & Ownership (\%) & $\Delta$ \\
\midrule
Yaari benchmark (SS on) & 100.0 & --- \\
+ Bequest motives (DFJ) & 100.0 & +0.0 pp \\
+ Minimum purchase (25K) & 74.1 & -25.9 pp \\
+ Pricing loads (MWR=0.85) & 59.5 & -14.7 pp \\
\midrule
\textit{Channels omitted by Pashchenko:} & & \\
+ Medical costs + R-S correlation & 38.2 & -21.3 pp \\
+ Survival pessimism & 13.7 & -24.4 pp \\
+ Full loads + Inflation & 8.5 & -5.3 pp \\
\midrule
Observed (Lockwood 2012) & 3.6 & --- \\
\bottomrule
\end{tabular}
\begin{tablenotes}
\small
\item Steps 0--3 incorporate the channels Pashchenko (2013) identified
(SS, bequests, minimum purchase, pricing loads) into our unified model.
Steps 4--6 add channels not in her framework. Note: this is not a
replication of her model (different preferences, health dynamics,
solution method). Housing illiquidity is not modeled; see text.
\end{tablenotes}
\end{table}



\subsection{Replication checks against prior literature}
\label{app:replication_checks}

The model implementation is checked against three prior papers in the annuity-puzzle literature. The checks are not full replications (the present model differs from each in nontrivial respects), but they verify that the implementation of shared components reproduces the published quantitative patterns.

\textbf{\citet{lockwood2012}.} The repository's automated test suite (\texttt{test/test\_lockwood.jl}) exercises eight checks against Lockwood's published moments: (i) the SSA cumulative death probability life table reproduces Lockwood's Table-A1 conditional survival at age 65 to 1e-8 tolerance; (ii) the fair payout rate at $r = 0.03$ matches Lockwood's reported 7.80\% annual payout per dollar of premium (his unit annuity price $\approx$12.82, i.e.\ $1/12.82 \approx 0.0780$; the model computes 0.07799); (iii) the willingness-to-pay (WTP) for a fair annuity without bequest motives reproduces Lockwood's $\approx$25\% of non-annuitized wealth; (iv) WTP declines monotonically with the pre-annuitized wealth fraction as Lockwood reports; (v) WTP collapses with bequest motives at the magnitudes Lockwood documents; (vi) loaded annuities reduce WTP consistent with Lockwood's loaded-pricing exhibits; (vii) the bequest-parameter calibration reproduces Lockwood's Table-1 estimates; (viii) the value function exhibits Lockwood's monotonicity and concavity properties. All checks pass under the present production grid.

\textbf{\citet{pashchenko2013}.} The Pashchenko comparison is reported as Table~\ref{tab:pashchenko}. Levels are not directly comparable (channel sets and preference specifications differ), but the qualitative finding is preserved: standard rational channels alone substantially overshoot observed US ownership. The channels Pashchenko omits (the combined Med+R-S correlation, survival pessimism, age-varying needs, and inflation erosion) reduce the residual but do not close it on their own. Adding the structural public-care aversion channel yields the nine-channel structural baseline of \ownNineChannelLTC{} (the disciplined reading).

\textbf{\citet{reichlingsmetters2015}.} Their core mechanism---the health-mortality correlation that makes annuities ``valuation-risky'' precisely when liquid wealth is most valuable---is operationalized in this paper as the combined Med+R-S channel (Section~\ref{M-sec:results}). The bundling reflects a structural property of the present framework: the R-S mechanism's quantitative bite operates through the interaction with stochastic medical costs, since without competing demand for liquid wealth in sick states the lower expected annuity NPV does not translate into a precautionary motive against annuitization. Activating the combined channel reduces predicted ownership by \magDeltaMedRS{}~pp in the sequential decomposition, contributing \shapNineMedRS{}~pp in the nine-channel structural Shapley attribution (\shapMedRS{}~pp in the eleven-channel exploratory game). The qualitative R-S sign-flip pattern (annuity demand falls when health-correlated medical risk is added on top of pure stochastic mortality) is reproduced by the present model under their parameter choices, verified in \texttt{test/test\_health.jl}.


%% ================================================================
%%  E. DIA/QLAC EXTENSION
%% ================================================================
\section{Deferred Income Annuity Extension}
\label{app:dia}

Deferred income annuities (DIAs) begin payments at a future age, typically 80 or 85. By concentrating payments in late life---when longevity risk is greatest---they offer higher per-dollar payouts. At the exhibit's nominal pricing, a DIA purchased at 65 with payments beginning at 80 pays roughly three times the per-period income of an immediate annuity (SPIA) per dollar of premium, and a DIA-85 roughly six times as much (Table~\ref{tab:dia} reports the payout rates).

Deferred products carry lower money's worth ratios, however. \citet{wettstein2021} estimate an MWR near 0.50 for deferred annuities beginning at age 85; we assume 0.50 for the DIA-80 and 0.45 for the DIA-85, extrapolating their deferral gradient, against \pMwrBaseline{} for SPIAs. The lower MWR reflects the longer deferral period (more discounting), smaller market (less competition), and higher sensitivity to mortality projection errors.

\begin{table}[htbp]
\centering
\caption{SPIA vs.\ Deferred Income Annuity (DIA) Comparison}
\label{tab:dia}
\begin{tabular}{lccc}
\toprule
 & SPIA & DIA-80 & DIA-85 \\
MWR & 0.87 & 0.50 & 0.45 \\
Payout rate & 0.0623 & 0.1659 & 0.3485 \\
\midrule
\multicolumn{4}{l}{\textit{No bequest}} \\
\quad Ownership (\%) & 1.6 & 0.0 & 0.4 \\
\quad CEV at \$200K (\%) & 0.00 & 0.00 & 0.00 \\
\quad Optimal $\alpha$ at \$200K & 0.0\% & 0.0\% & 0.0\% \\
\midrule
\multicolumn{4}{l}{\textit{Moderate (DFJ)}} \\
\quad Ownership (\%) & 0.0 & 0.0 & 0.0 \\
\quad CEV at \$200K (\%) & 0.00 & 0.00 & 0.00 \\
\quad Optimal $\alpha$ at \$200K & 0.0\% & 0.0\% & 0.0\% \\
\bottomrule
\end{tabular}
\begin{tablenotes}
\small
\item DIA MWR from Wettstein et al.\ (2021). All models include
medical costs, R-S correlation, pricing loads, and 2\% inflation.
CEV evaluated at good health (H=1).
\end{tablenotes}
\end{table}


Table~\ref{tab:dia} compares predicted ownership and welfare across product types. The exhibit is solved without pre-existing Social Security or defined-benefit income, so its SPIA ownership levels are not comparable to the headline configuration; the comparison is internal to the table's common environment. The DIA results suggest that product innovation alone is unlikely to resolve low annuity demand. The channels suppressing SPIA demand---valuation risk, survival pessimism, inflation erosion, bequest motives---also bear on deferred products, and the inferior pricing of DIAs partially offsets their actuarial advantage.


%% ================================================================
%%  F. FULL ROBUSTNESS TABLE
%% ================================================================
\section{Full Sensitivity Analysis}
\label{app:robustness}

Tables~\ref{tab:full_robustness} and~\ref{tab:full_robustness_continued} report predicted ownership under all robustness specifications.

\begin{table}[htbp]
\centering
\caption{Full Sensitivity Results}
\label{tab:full_robustness}
\begin{threeparttable}
\small
\begin{tabular}{llc}
\toprule
Category & Specification & Ownership (\%) \\
\midrule
\multicolumn{3}{l}{\textit{Risk aversion ($\gamma$)}} \\
 & $\gamma = 1.50$ & \ownGammaOnePointFiveNum \\
 & $\gamma = 2.00$ & \ownGammaTwoNum \\
 & $\gamma = 2.20$ & \ownGammaTwoPointTwoNum \\
 & $\gamma = 2.30$ & \ownGammaTwoPointThreeNum \\
 & $\gamma = 2.40$ & \ownGammaTwoPointFourNum \\
 & $\gamma = \pGamma$ (baseline) & \ownGammaTwoPointFiveNum \\
 & $\gamma = 2.60$ & \ownGammaTwoPointSixNum \\
 & $\gamma = 3.00$ & \ownGammaThreeNum \\
 & $\gamma = 3.50$ & \ownGammaThreePointFiveNum \\
 & $\gamma = 4.00$ & \ownGammaFourNum \\
 & $\gamma = 5.00$ & \ownGammaFiveNum \\[4pt]
\multicolumn{3}{l}{\textit{Money's worth ratio}} \\
 & MWR = 0.82               & \ownMWREightyTwoNum \\  % ALLOWLIST: sensitivity-table row showing prediction at a 0.82 load
 & MWR = 0.85               & \ownMWREightyFiveNum \\
 & MWR = \pMwrBaseline{} (baseline) & \ownNineChannelFullNum \\
 & MWR = 0.90               & \ownMWRNinetyNum \\
 & MWR = 0.95               & \ownMWRNinetyFiveNum \\[4pt]
\multicolumn{3}{l}{\textit{Hazard multipliers $[\mu_G, \mu_F, \mu_P]$}} \\
 & R--S functional $[0.45, 1.0, 3.5]$ & \ownHazardRSNum \\
 & Baseline $[\pHazardGood, \pHazardFair, \pHazardPoor]$ & \ownNineChannelFullNum \\
 & HRS SRH $[0.57, 1.0, 2.7]$ & \ownHazardHRSNum \\
 & Conservative $[0.60, 1.0, 2.0]$ & \ownHazardConservativeNum \\[4pt]
\multicolumn{3}{l}{\textit{Inflation ($\pi$)}} \\
 & $\pi = 1\%$ & \ownInflationOneNum \\
 & $\pi = \pInflation$ (baseline) & \ownInflationTwoNum \\
 & $\pi = 3\%$ & \ownInflationThreeNum \\[4pt]
\multicolumn{3}{l}{\textit{Survival pessimism ($\psi$)}} \\
 & $\psi = 0.960$ (strong-pessimism endpoint) & \ownPsiEndpointNum \\
 & $\psi = 0.970$ & \ownPsiNinetySevenNum \\
 & $\psi = \pPessimism$ (baseline) & \ownPsiNinetyEightOneNum \\
 & $\psi = 0.990$ & \ownPsiNinetyNineNum \\
 & $\psi = 1.000$ (objective) & \ownPsiObjectiveNum \\
\bottomrule
\end{tabular}
\begin{tablenotes}
\small
\item Unless otherwise noted, all specifications use baseline parameters: $\gamma = \pGamma$, $\beta = \pBeta$, DFJ bequests ($\theta = \pThetaDFJ$, $\kappa = \pKappaDFJ$), survival pessimism $\psi = \pPessimism$, MWR $= \pMwrBaseline$, $F = \pFixedCost$, $\pi = \pInflation$, hazard multipliers $[\pHazardGood, \pHazardFair, \pHazardPoor]$, \pNQuad-node Gauss--Hermite quadrature, production grid ($\pNWealth \times \pNAnnuity \times \pNAlpha$). All sweeps are reported on the nine-channel structural baseline (\ownNineChannelLTC{}, the disciplined reading).
\end{tablenotes}
\end{threeparttable}
\end{table}

\begin{table}[htbp]
\centering
\caption{Full Sensitivity Results (continued)}
\label{tab:full_robustness_continued}
\begin{threeparttable}
\small
\begin{tabular}{llc}
\toprule
Category & Specification & Ownership (\%) \\
\midrule
\multicolumn{3}{l}{\textit{Bequest specification}} \\
 & No bequest ($\theta = 0$) & \ownBequestNoneNum \\
 & DFJ luxury (baseline) & \ownNineChannelFullNum \\[4pt]
\multicolumn{3}{l}{\textit{Minimum purchase}} \\
 & \$0 & \ownNineChannelFullNum \\
 & \$1,000 & \ownNineChannelFullNum \\
 & \$5,000 & \ownNineChannelFullNum \\
 & \pMinPurchase{} (baseline) & \ownNineChannelFullNum \\
 & \$25,000 & \ownNineChannelFullNum \\
\bottomrule
\end{tabular}
\begin{tablenotes}
\small
\item Notes as in Table~\ref{tab:full_robustness}. Predicted ownership is invariant to the minimum-purchase threshold because model owners are concentrated at high wealth and purchase annuity premiums well above \$25{,}000, so the floor never binds (in contrast to the binding \$25{,}000 minimum in the Pashchenko comparison, Table~\ref{tab:pashchenko}). Grid and quadrature convergence are reported separately in Table~\ref{tab:grid_convergence}.
\end{tablenotes}
\end{threeparttable}
\end{table}


\subsection{Robustness to channel partition}
\label{app:partition}

The nine-channel game combines medical-expenditure risk with the Reichling--Smetters correlation into one channel, and combines Social Security and defined-benefit pension income into one pre-existing-annuitization channel. Both groupings reflect the economics: the R-S correlation has quantitative bite in this framework only through its interaction with medical risk, and SS and DB income enter the budget constraint identically, as non-commutable real annuity flows. To confirm that the channel \emph{ranking} does not depend on these choices, Table~\ref{tab:partition_robustness} recomputes the exact Shapley values over two alternative ten-channel partitions: Panel A unbundles medical risk from the R-S correlation, and Panel B splits Social Security income from DB pension income. In both partitions the ranking is preserved---pricing loads remain the dominant suppressor of demand and bequest motives remain mid-pack---so the headline ordering is an artifact of neither bundling choice.

\input{../tables/tex/partition_robustness.tex}

%% ================================================================
%%  G. GAUSS-HERMITE QUADRATURE CHECK
%% ================================================================
\section{Gauss--Hermite Quadrature Check}
\label{app:quadrature}

Medical expenditure shocks are integrated via Gauss--Hermite quadrature over the lognormal distribution $\ln m \sim N(\mu, \sigma^2)$ with $\sigma \approx 1.4$. At this variance, the distribution has heavy tails that interact with the Medicaid consumption floor to create a kink in the integrand. Low-order quadrature rules place insufficient weight in the tails, producing inaccurate expectations.

Table~\ref{tab:grid_convergence} (Panel B) documents the convergence: binary ownership at the production grid moves from \quadThree{} at $n_{\text{quad}} = 3$ to \quadNine{} at $n_{\text{quad}} = 9$ and rises within a \quadBand{} pp window thereafter (\quadNine, \quadEleven, \quadThirteen, \quadFifteen{} at $n_{\text{quad}} \in \{9, 11, 13, 15\}$). At the coarsest configurations the level scatter is wider than the refinement path alone suggests (the 40$\times$15 grid gives \gridConvCoarse{}), reinforcing the paper's practice of quoting the level as a range. Lower-order rules deviate by up to \quadLowDeviation{} pp because of the tail--floor interaction: small changes in the quadrature approximation of the medical-expense tail shift marginal households across the discrete ownership threshold. Nine nodes balances accuracy against computational cost; we adopt it as the production standard.

For reference, the exact lognormal moment $E[\exp(\sigma \varepsilon)] = \exp(\sigma^2/2)$ is matched to machine precision by all rules with $n \geq 5$. The residual quadrature error operates through the nonlinear interaction of the medical expense realization with the consumption policy function and Medicaid floor, not through moment inaccuracy.


%% ================================================================
%%  PAIRWISE CHANNEL INTERACTIONS
%% ================================================================
\section{Pairwise Channel Interactions}
\label{app:pairwise}

Table~\ref{tab:pairwise} reports the pairwise interaction strength for the eight rational and preference channels. For each pair (A, B), the interaction is defined as the difference between the ownership drop when both are active simultaneously and the sum of their individual drops. A negative interaction indicates super-additivity: the channels reinforce each other. The order-independent contributions of the nine structural channels are reported in the Shapley decomposition (Section~\ref{M-sec:shapley}); the exploratory eleven-channel extension is in Appendix~\ref{app:behavioral_sensitivity}.

\begin{table}[htbp]
\centering
\caption{Pairwise Interaction Strengths, Rational and Preference Channels (pp)}
\label{tab:pairwise}
\begin{tabular}{lcccccccc}
\toprule
  & SS & Bequests & Medical+R-S & Loads & Inflation & Pessimism & HealthUtil & AgeNeeds \\
\midrule
SS & --- & +0.0 & -1.3 & -31.3 & +2.3 & -11.5 & +0.2 & -0.4 \\
Bequests & --- & --- & -0.3 & -0.1 & -0.0 & -0.1 & +0.0 & -0.1 \\
Medical+R-S & --- & --- & --- & -2.5 & +0.8 & -0.6 & +0.4 & -0.6 \\
Loads & --- & --- & --- & --- & -0.1 & -0.3 & +0.0 & +0.0 \\
Inflation & --- & --- & --- & --- & --- & +0.9 & +0.0 & +0.4 \\
Pessimism & --- & --- & --- & --- & --- & --- & +0.2 & +0.1 \\
HealthUtil & --- & --- & --- & --- & --- & --- & --- & -0.1 \\
AgeNeeds & --- & --- & --- & --- & --- & --- & --- & --- \\
\bottomrule
\end{tabular}
\begin{tablenotes}
\small
\item Each cell shows the interaction: ownership with both channels minus
the additive prediction from individual effects. Negative values indicate
super-additive demand reduction (channels reinforce each other).
\end{tablenotes}
\end{table}


The strongest interaction is between Social Security and pricing loads ($\pairwiseSSLoads$ pp). SS raises demand by providing an income floor, and loads destroy demand by eliminating the surplus; when combined, the effect of loads is amplified by the elevated demand that SS creates. The combined Med+R-S health-mortality channel and pricing loads interact at $\pairwiseMedRSLoads$ pp, reflecting the economic complementarity between valuation risk and pricing frictions. Bequest interactions are near zero across the board, consistent with the modest isolated effect of the DFJ luxury-good specification.


%% ================================================================
%%  SUFFICIENT-STATISTIC FACTORIZATION (HEURISTIC)
%% ================================================================
\section{Sufficient-Statistic Factorization (Heuristic)}
\label{app:phi_factorization}

Each demand-suppressing channel introduces a factor $\Phi_i < 1$ that reduces the annuity value ratio $R$ of Section~\ref{M-sec:sufficient_stats} below the Yaari benchmark. As an approximate sufficient-statistic representation (exact only under restrictions stronger than the finite-horizon stochastic problem solved here), the annuity value ratio can be written as a product of channel-specific factors,
\begin{equation}
\label{eq:multiplicative_R}
R \approx \Phi_{\text{MWR}} \times \Phi_{\text{bequest}} \times \Phi_{\text{health}} \times \Phi_{\text{inflation}} \times \Phi_{\text{SS}} \times \Phi_{\text{pessimism}} \times \Phi_{\text{age}} \times \Phi_{\text{state}},
\end{equation}
where:
\begin{itemize}
\item $\Phi_{\text{MWR}} = \text{MWR}$. The pricing load directly scales the numerator. At MWR $= \pMwrBaseline$, each dollar of premium buys \pMwrBaseline{} of expected value, a load of \pMwrLoad{} relative to the actuarially fair benchmark.

\item $\Phi_{\text{bequest}} = \big[1 + (1 - s) \, v'(b) / u'(c)\big]^{-1}$, where $s$ is the survival probability and $v'(b)/u'(c)$ is the ratio of marginal bequest to marginal consumption utility. Under the DFJ luxury-good specification ($\kappa = \pKappaDFJ \gg$ median wealth $\approx \$40{,}000$), $v'(b) \approx \theta(\kappa)^{-\gamma}$ is nearly constant and small, so $\Phi_{\text{bequest}} \approx 1.00$.

\item $\Phi_{\text{health}}$ reflects the covariance between survival and the marginal utility of liquid wealth. When health deterioration simultaneously reduces survival and raises medical costs, the ratio $V_A / V_W$ falls. This is the Reichling--Smetters mechanism.

\item $\Phi_{\text{inflation}}$ equals the ratio of the present value of the nominal payment stream (in real terms) to the real payment stream, weighted by marginal utilities. At 2\% inflation, payments in year 25 retain only 61\% of their initial purchasing power.

\item $\Phi_{\text{SS}}$ captures the diminishing marginal insurance value of additional annuitization when Social Security already provides a longevity-insured income floor.

\item $\Phi_{\text{pessimism}} < 1$ because the agent uses $\psi \cdot s(t,H)$ in the Bellman equation. The subjective present value of future annuity income is lower than its objective present value.

\item $\Phi_{\text{age}} < 1$ because declining consumption needs $w(t)$ reduce the marginal utility of late-life income, which is precisely the income the annuity provides.

\item $\Phi_{\text{state}} \leq 1$ because poor-health states generate lower marginal utility of consumption, and the annuity pays into those states.
\end{itemize}
This factorization is a heuristic guide to how the channels interact, not a proven identity; the exact, model-consistent attribution is the numerical Shapley decomposition of Section~\ref{M-sec:shapley}.

The decomposition in equation~(\ref{eq:multiplicative_R}) characterizes the nine-channel structural baseline. SDU enters the value function as a separate multiplicative factor on portfolio-financed consumption (equation~\ref{M-eq:sdu}), and the at-purchase penalty (equation~\ref{M-eq:ped}) enters as a one-time additive shifter on the annuitization decision; together they contribute to the eleven-channel Shapley enumeration in Section~\ref{M-sec:shapley}.

\paragraph{Super-additivity of participation effects (heuristic).} In the same approximate representation, the fraction of agents annuitizing is $\Pr(\alpha^* > 0) = \Pr(R_i > 1)$, where $R_i$ varies across agents due to heterogeneity in wealth, health, and pre-existing annuitized income. In log space,
\begin{equation}
\ln R_i \approx \sum_j \ln \Phi_j + \varepsilon_i,
\end{equation}
where $\varepsilon_i$ captures agent-level heterogeneity. Each additional channel shifts the distribution of $\ln R_i$ leftward by $|\ln \Phi_j|$. Agents near the participation threshold ($R_i \approx 1$, i.e., $\ln R_i \approx 0$) are disproportionately eliminated, so the ownership drop from adding a channel is larger when other channels have already thinned the right tail of the $R_i$ distribution. This super-additive structure motivates the exact Shapley decomposition in Section~\ref{M-sec:shapley}, which attributes the total demand reduction to individual channels without dependence on the order in which channels are added.

The dominant term is $\Phi_{\text{MWR}} = \pMwrBaseline$, which alone reduces the annuity value ratio by \pMwrLoad. The near-unity $\Phi_{\text{bequest}}$ under the DFJ specification explains why bequests contribute only \shapNineBequests{} pp in the Shapley decomposition: the luxury-good curvature parameter makes bequest utility irrelevant for median-wealth retirees.


%% ================================================================
%%  H. PERIOD-CERTAIN ANNUITY
%% ================================================================
\section{Period-Certain Annuity}
\label{app:period_certain}

A 10-year period-certain annuity guarantees payments for the first 10 years regardless of survival: in market contracts, death during the guarantee window passes the remaining guaranteed payments (or their present value) to the estate, partially offsetting the bequest penalty of life-only annuities.

The period-certain payout rate is lower than the life-only rate because the insurer cannot reclaim payments from early decedents during the guarantee period. The pricing formula replaces survival probabilities with 1.0 for the first 10 years:
\[
\text{PV}_{\text{certain}} = \sum_{k=0}^{9} \frac{1}{(1+r_{\text{nom}})^k} + \sum_{k=10}^{T-1} \frac{\prod_{j=0}^{k-1} s_{\text{base}}(j)}{(1+r_{\text{nom}})^k}.
\]
Under the production mortality table and discounting, the 10-year guarantee lowers the fair payout rate by \pctPeriodCertainCutNom{} under nominal pricing (\pctPeriodCertainCutReal{} under real pricing); the forgone first-decade mortality credit funds the guarantee.

Only the pricing side of this product is implemented in the replication package (\texttt{emit\_\allowbreak period\_\allowbreak certain\_\allowbreak pricing\allowbreak .jl}). The household side---crediting the estate with the present value of remaining guaranteed payments at death---does not enter the value function, so the model does not compute the net demand effect of the guarantee, which is ambiguous a priori: the guarantee reduces the bequest penalty but gives up part of the mortality credit that makes annuitization attractive in the first place. Two features of the calibrated model bound the plausible upside: bequest motives contribute \shapNineBequests{} pp in the nine-channel Shapley decomposition---mid-pack---and under the DFJ luxury-good specification the marginal bequest utility of median-wealth retirees is low, so a product margin that operates only through the bequest penalty is unlikely to overturn the loads-dominated ranking. A full household-side valuation of guarantee riders is left as an extension.


%% ================================================================
%%  I. IMPLIED RISK AVERSION
%% ================================================================
\section{Joint Parameter Uncertainty: Draw Design}
\label{app:mc_design}

The joint stress test in Section~\ref{M-sec:joint_uncertainty} draws \nMCDraws{} parameter vectors from independent uniform distributions: the Poor-health hazard multiplier $\mu_P \sim U(2.5, 5.0)$, inflation $\pi \sim U(0.015, 0.025)$, MWR $\sim U(0.83, 0.91)$, survival beliefs $\psi \sim U(0.92, 1.00)$ (wider than the reported $[0.96, 1.0]$ calibration range, extending the stress test mildly beyond the strong-pessimism endpoint), and the consumption-needs slope $\delta_c \sim U(0.01, 0.03)$. Risk aversion is held fixed at the baseline. Each of the other ranges is symmetric around its production value (e.g., $\mu_P$ around $\pHazardPoor$), so the central calibration is its central tendency; the $\psi$ draw is centered on the strong-pessimism endpoint (0.96) and covers both the baseline $\pPessimism$ and objective beliefs. For each draw the model is solved on a convergence-verified $60 \times 20 \times 51$ grid, coarser than production for tractability across \nMCDraws{} draws.

The implied-$\gamma$ exercise of Section~\ref{app:monte_carlo} uses different ranges by design: it draws the Poor-health multiplier over the narrower empirical anchor range $U(2.0, 3.5)$ (HRS self-reported to R-S functional) and MWR over the wider $U(0.75, 0.89)$, whose lower reach spans the classic population-table estimates.

\section{Implied Risk Aversion from Parameter Uncertainty}
\label{app:monte_carlo}

The decomposition results depend on calibrated parameters, and the sensitivity to $\gamma$ documented in Section~\ref{M-sec:monte_carlo} raises the question of whether the baseline prediction is fragile. To address this, 300 draws of nuisance parameters are taken from the following distributions:

\begin{itemize}
\item \textbf{Poor-health hazard multiplier:} $\mu_P \sim U(2.0, 3.5)$. The lower bound corresponds to conservative SRH-based estimates for the oldest old. The upper bound corresponds to the functional limitation states used by \citet{reichlingsmetters2015}. The Good-health multiplier is held at 0.50.

\item \textbf{Inflation rate:} $\pi \sim U(0.01, 0.03)$. This spans near-term CPI uncertainty around the baseline calibration of 2\%.

\item \textbf{Money's worth ratio:} $\text{MWR} \sim U(0.75, 0.89)$. This reflects measurement uncertainty in current annuity pricing estimates from \citet{mitchell1999} and \citet{wettstein2021}; the production baseline is $\pMwrBaseline$.
\end{itemize}

For each draw, the model is solved repeatedly, bisecting over $\gamma$ to target Lockwood's 3.6\% observed ownership; the search terminates when the $\gamma$ bracket narrows below 0.01. Table~\ref{tab:implied_gamma} reports the distribution of implied $\gamma$ values.

\begin{table}[htbp]
\centering
\caption{Implied Risk Aversion from Monte Carlo Parameter Uncertainty}
\label{tab:implied_gamma}
\begin{tabular}{lc}
\toprule
Statistic & Value \\
\midrule
Number of draws & 300 \\
Target ownership & 3.6\% \\
Median implied $\gamma$ & 2.45 \\
Mean implied $\gamma$ & 2.48 \\
Interquartile range & [2.37, 2.57] \\
Min / Max & 2.19 / 2.88 \\
Fraction in [1.5, 3.0] & 100\% \\
\bottomrule
\end{tabular}
\begin{tablenotes}
\small
\item For each draw of nuisance parameters ($\mu_P$, $\pi$, MWR),
we bisect over $\gamma$ to find the value that generates 3.6\%
predicted ownership (Lockwood 2012 observed rate).
Draws: $\mu_P \sim U(2.0, 3.5)$, $\pi \sim U(0.01, 0.03)$, MWR $\sim U(0.75, 0.89)$.
Survival pessimism $\psi = 0.981$ (O'Dea \& Sturrock 2023).
\end{tablenotes}
\end{table}


The median implied $\gamma$ sits inside the \citet{chetty2006} range of 1--3, with all 300 converged draws falling inside the Chetty interval and a tight interquartile range (Table~\ref{tab:implied_gamma}). Matching Lockwood's reported 3.6\% rate therefore requires a $\gamma$ close to the baseline $\gamma = \pGamma$ rather than an implausible value. This is reassuring about the channel structure adopted in the body of the paper: the baseline risk aversion is consistent with the ownership rate Lockwood reports, without recourse to extreme preferences. Adding the two preference channels brings the baseline-$\gamma$ prediction to \ownEightChannelExt{} and the structural public-care aversion channel yields the nine-channel structural baseline of \ownNineChannelLTC{} (the disciplined reading).

For each draw, the full model is solved by backward induction on a reduced $60 \times 20 \times 51$ grid---coarser than the production $\pNWealth \times \pNAnnuity \times \pNAlpha$ grid, to keep the 300-draw bisection exercise tractable---and ownership is evaluated on the HRS population sample. The 300 bisection searches (each requiring 8--12 solves) are parallelized across available CPU cores.


%% ================================================================
%%  J. SOLUTION ALGORITHM DETAILS
%% ================================================================
\section{Solution Algorithm}
\label{app:algorithm}

\subsection{Backward induction}

The model is solved by backward induction from the terminal period $T = 46$ (age 110) to $t = 1$ (age 65). At the terminal period, survival probability is zero and the individual divides available resources between consumption and a terminal bequest:
\[
V(W, A, H, T) = \max_{c \leq W + A_T + \text{SS}(T) - m_T} \Big\{ u(c) + \beta \, v\big(W + A_T + \text{SS}(T) - m_T - c\big) \Big\}.
\]

For $t = T-1, \ldots, 1$, the value function is computed at each grid point $(W_i, A_j, H_k)$ by solving the one-period consumption problem using Brent's method on the interval $[\underline{c}, \; W + A_t + \text{SS}(t) - m]$. The continuation value $V(W', A', H', t+1)$ is evaluated by piecewise linear interpolation over the wealth grid with flat extrapolation beyond the grid boundaries.

\subsection{Expectation computation}

The expectation over next-period health and medical expenditures involves:
\begin{enumerate}
\item A sum over three health states weighted by transition probabilities $\pi(H_t, H', t)$.
\item A sum over nine Gauss--Hermite quadrature nodes for the medical expenditure shock $\varepsilon$.
\end{enumerate}
Each grid point requires $3 \times 9 = 27$ evaluations of the continuation value, yielding a total of $80 \times 30 \times 3 \times 46 \times 27 \approx 8.94$ million value function evaluations per model solve.

The optimal consumption policy is computed by integrating over the joint distribution of next-period health and medical expenditure shocks; the object carried forward is the value function $V(\cdot, t+1)$. Forward simulations are shock-contingent: at each simulated period the consumption choice is re-optimized against the realized health state and medical-expenditure draw using the stored continuation value, rather than applying a precomputed certainty-equivalent policy. This matters because the medical-expenditure distribution is heavy-tailed and interacts with the Medicaid consumption floor, so applying an averaged policy to realized tail draws would misstate spend-down behavior near the floor.

\subsection{Annuitization optimization}

At $t = 1$, the optimal annuitization fraction $\alpha^*$ is found by grid search over 101 equally spaced values in $[0, 1]$. For each candidate $\alpha$, post-purchase wealth is $(1 - \alpha)W_0$ and annuity income is $\alpha W_0 \times \text{payout\_rate}$. The value function at the resulting state is evaluated by bilinear interpolation over the $(W, A)$ grid for the relevant health state. The fixed cost $F$ is subtracted from post-purchase wealth whenever $\alpha > 0$.

\subsection{Population evaluation}

Ownership is evaluated on a population of HRS individuals with observed initial wealth, age, and self-reported health. Pre-existing annuitized income $y$ is assigned at the wealth-band level (the four $\text{SS}_{\text{obs}} + \text{DB}_{\text{obs}}$ levels of Table~\ref{M-tab:calibration}), not per individual; for each individual, the optimal $\alpha^*$ is then computed at their specific $(W_0, y_{b(W_0)}, H_0)$. Ownership is defined as $\alpha^* > 0$. The ownership rate is the fraction of the eligible population (wealth $\geq$ \pMinWealth) that annuitizes any positive amount. The wealth grid extends to \pWMaxMillions, covering the 99.5th percentile of the HRS wealth distribution. Individuals above the grid maximum (\pctHRSAboveWmax{} of the sample) are capped at the grid maximum and evaluated at the cap, rather than extrapolated beyond it. A related bound applies to the annuity-income grid, which tops out at the full-annuitization payout of the grid-maximum wealth: for 9 of the 2{,}279 eligible observations (0.4\%) the implied annuity income at $\alpha = 1$ exceeds the grid ceiling, by at most 5.9\% (values audited by \texttt{test/test\_grid\_clamp\_audit.jl}) (the mechanism is the nominal premium gross-up at later purchase ages, $(1+\pi)^{t-1}$, rather than pre-existing income, which enters through the income function rather than the annuity state), which biases against high-purchase allocations and is therefore conservative for the low-ownership findings.

\subsection{Computational performance}

A single model solve (one parameterization) requires approximately 190 seconds on an Apple M1 processor with 9-node Gauss--Hermite quadrature. The full decomposition (8 steps with per-quartile SS separation) runs in approximately 25 minutes. The robustness analysis (approximately 40 configurations using mean SS) requires 2--3 hours. The Monte Carlo exercise (1,000 solves) requires approximately 50 hours with parallelization.


%% ================================================================
%%  K. EULER EQUATION RESIDUALS
%% ================================================================
\section{Euler Equation Residuals}
\label{app:euler}

The intertemporal optimality condition at each interior grid point $(W, A, H, t)$ is
\[
u'(c^*) = \beta(1+r) \, \mathbb{E}_m\!\left[ s(t,H) \sum_{H'} \pi(H,H',t) \, V_W(W', A, H', t{+}1) + (1{-}s(t,H)) \, v'(W') \right],
\]
where $W' = (1+r)(\text{cash} - c^*)$ and the expectation is over the medical expense shock $m$. The marginal value of next-period wealth $V_W$ is computed by finite differences on the piecewise-linear value function interpolant.

The residuals are computed for the structural configuration (source-dependent utility off); the SDU channel's kinked marginal utility is not covered by this diagnostic. The normalized Euler residual at each grid point is $|u'(c^*) - \text{RHS}| / |u'(c^*)|$. Corner solutions ($c^* = \underline{c}$ or $c^* = \text{cash}$) are excluded, as the Euler equation holds with inequality at boundaries.

Table~\ref{tab:euler} reports summary statistics. The median residual is below $10^{-6}$; the summary statistics are computed over strictly positive residuals, so no grid point satisfies the condition to machine zero, but half sit below that threshold. The mean residual is \eulerMeanPct, with approximately \eulerPctAboveOne{} of grid points exceeding 1\% and \eulerPctAboveFive{} exceeding 5\%. The maximum residual stays below 1 at the production grid and rises slightly above 1 at the finer $100 \times 40$ grid; these large values are confined to low-wealth states where the finite-difference derivative of the piecewise-linear value function poorly approximates the true marginal value, and the maximum does not fall monotonically with refinement. The share of points above 1\% is stable across resolutions, easing from \eulerPctOneCoarse{} at $40 \times 15$ to \eulerPctOneFine{} at $100 \times 40$.

\input{../tables/tex/euler_residuals_table.tex}

The use of piecewise-linear interpolation means derivatives are discontinuous at grid points, which introduces a floor on residual accuracy. The \eulerPctAboveFive{} of points with residuals above 5\% are concentrated at low wealth levels where the value function has high curvature and the linear approximation is poorest. Switching to cubic spline interpolation would reduce these residuals at the cost of potential Runge-type oscillations at low wealth. The grid convergence results in Section~\ref{app:grid_convergence} provide a complementary accuracy check: stable ownership predictions across grid resolutions confirm that the policy function is resolved by the production grid. The exact Shapley decomposition, which evaluates all $2^{9} = 512$ subsets of the nine-channel structural game (and, in the exploratory extension, all $2{,}048$ subsets of the eleven-channel game) on the same production grid, implicitly provides a further robustness check: if the solver were unreliable, these ownership evaluations would not produce the internally consistent attribution hierarchy reported in Section~\ref{M-sec:shapley}.


\bibliographystyle{aer}
%% ================================================================
%%  LOADS-SPLIT SHAPLEY DECOMPOSITION
%% ================================================================
\section{Loads-Split Shapley Decomposition}
\label{app:loads_split}

The nine-channel game bundles the money's worth wedge, the fixed purchase cost, and the minimum-purchase requirement as one pricing-loads player. Table~\ref{tab:shapley_loads_split} reruns the game with the three components as separate players (eleven players, $2^{11} = 2{,}048$ subsets, production grid); the full coalition reproduces the nine-channel full model exactly (\splitFullOwn). The money's worth wedge remains the single largest suppressor (\splitMWRWedge~pp), narrowly ahead of the fixed purchase cost (\splitFixedCost~pp), the next-largest suppressor. The fixed cost itself exceeds bequest motives (\splitBequests~pp), and the minimum-purchase floor is negligible (\splitMinPurchase~pp), consistent with the invariance diagnostic in Appendix~\ref{app:robustness}. Unbundling narrows the wedge's margin over the second tier relative to the bundled channel's value; Shapley values are not additive across game partitions, so the three components need not (and do not) sum to the bundled attribution. The finer partition preserves the claim in the title: the proportional pricing wedge alone out-ranks every other structural channel, and each of the two larger loads components individually exceeds the bequest attribution.

\input{../tables/tex/shapley_loads_split.tex}

%% ================================================================
%%  SEX-BLENDED MORTALITY ROBUSTNESS
%% ================================================================
\section{Mortality-Convention Robustness}
\label{app:male_mortality}

The headline calibration uses the sex-blended, aggregate-anchored life table (Section~\ref{M-sec:calibration}). Table~\ref{tab:shapley_male} recomputes the full nine-channel game under the prior convention in this literature: the male SSA 2003 table from the \citet{lockwood2012} replication files, with unnormalized health hazards, used for both survival and pricing. The ranking is unchanged: pricing loads contribute \maleTableLoads~pp and bequest motives remain mid-pack (\maleTableBequests~pp). The predicted level is \maleTableOwn{} under the prior convention, against the headline level under the blended anchored table---the level moves with the mortality treatment while the ranking does not.

\input{../tables/tex/shapley_male_mortality.tex}

%% ================================================================
%%  ADDITIONAL TABLES AND FIGURES
%% ================================================================
\section{Additional Tables and Figures}
\label{app:additional_exhibits}

Display items referenced from the main text: the joint parameter-uncertainty summary, the lifecycle moment check, the cross-sectional ownership logit, the Social Security cut schedule, the welfare counterfactual CEVs, and the risk-aversion, hazard-multiplier, policy-function, and welfare-surface figures.

\begin{table}[htbp]
\centering
\caption{Conditional Monte Carlo: Predicted Ownership at $\gamma = 2.5$}
\label{tab:monte_carlo}
\begin{tabular}{lc}
\toprule
Statistic & Value \\
\midrule
Number of draws & 1000 \\
Median predicted ownership & 31.9\% \\
90\% sensitivity interval & [17.8\%, 43.3\%] \\
Interquartile range (50\% interval) & [26.5\%, 36.5\%] \\
Mean & 31.4\% \\
Min / Max & 7.0\% / 53.2\% \\
Fraction in [1\%, 10\%] & 0\% \\
Fraction in [3\%, 6\%] (observed range) & 0\% \\
\bottomrule
\end{tabular}
\begin{tablenotes}
\small
\item Risk aversion fixed at $\gamma = 2.5$. Joint draws over six
calibration-uncertain parameters: $\mu_P \sim U(2.0, 3.5)$, $\pi \sim U(0.015, 0.025)$,
MWR $\sim U(0.83, 0.91)$, $\psi \sim U(0.97, 1.0)$, $\delta_c \sim U(0.01, 0.03)$,
$\psi_{\text{purchase}} \sim U(0.005, 0.030)$. Full ten-channel model.
\end{tablenotes}
\end{table}

\begin{table}[htbp]
\centering
\caption{Simulated vs Empirical Lifecycle Moments}
\label{tab:moment_validation}
\begin{tabular}{lcc}
\toprule
Moment & Simulated & Empirical (HRS) \\
\midrule
Mean bequest & \$40023 & \$90000 \\
Median bequest & \$0 & \$20000 \\
Fraction bequest > \$10K & 37.0\% & 45.0\% \\
\bottomrule
\end{tabular}
\begin{tablenotes}
\small
\item Simulated: 100,000 trajectories, initial wealth \$250,000, Fair health.
\item Empirical: HRS exit interviews (bequests); Jones et al.\ (2018) (medical).
\end{tablenotes}
\end{table}

\input{../tables/tex/empirical_gradients_logit.tex}
\begin{table}[htbp]
\centering
\caption{Private Annuity Demand Response to Social Security Benefit Reductions}
\label{tab:ss_cut}
\begin{tabular}{lccc}
\toprule
SS Benefit Cut & Ownership (\%) & Mean $\alpha$ & $\Delta$ (pp) \\
\midrule
0\% (baseline) & 1.8 & 0.002 & --- \\
10\% & 3.8 & 0.005 & +2.1 \\
15\% & 5.9 & 0.007 & +4.1 \\
23\% (trust fund) & 10.4 & 0.015 & +8.7 \\
30\% & 16.3 & 0.030 & +14.5 \\
40\% & 25.1 & 0.059 & +23.3 \\
50\% & 28.6 & 0.084 & +26.8 \\
100\% (elimination) & 14.3 & 0.067 & +12.5 \\
\midrule
Observed (Lockwood 2012) & 3.6 & & \\
\bottomrule
\end{tabular}
\begin{tablenotes}
\small
\item Full model with all channels active. SS quartile levels
scaled by (1 -- cut fraction). Baseline levels: [\$14K, \$17K, \$20K, \$25K].
23\% cut corresponds to projected trust fund exhaustion circa 2033.
100\% cut is a theoretical benchmark (complete SS elimination).
\end{tablenotes}
\end{table}

\begin{table}[htbp]
\centering
\caption{Welfare Gain from Annuity Access Under Policy Counterfactuals (CEV, \%)}
\label{tab:cev_counterfactuals}
\begin{tabular}{lcccc}
\toprule
Wealth & Baseline & Group pricing & Real annuity & Best feasible \\
\midrule
\multicolumn{5}{l}{\textit{Panel A: Good Health, DFJ Bequests}} \\
\$50K & 0.00 & 0.00 & 0.00 & 0.00 \\
\$100K & 0.00 & 0.00 & 0.00 & 0.00 \\
\$200K & 0.00 & 0.00 & 0.00 & 1.47 \\
\$500K & 0.24 & 0.75 & 0.42 & 4.43 \\
\$1000K & 0.86 & 1.53 & 1.07 & 5.52 \\
\midrule
\multicolumn{5}{l}{\textit{Panel B: Fair Health, DFJ Bequests}} \\
\$50K & 0.00 & 0.00 & 0.00 & 0.00 \\
\$100K & 0.00 & 0.00 & 0.00 & 0.00 \\
\$200K & 0.00 & 0.00 & 0.00 & 0.80 \\
\$500K & 0.07 & 0.49 & 0.21 & 3.64 \\
\$1000K & 0.59 & 1.12 & 0.77 & 4.59 \\
\midrule
\multicolumn{5}{l}{\textit{Panel C: Good Health, No Bequests}} \\
\$50K & 0.00 & 0.00 & 0.00 & 0.00 \\
\$100K & 0.00 & 0.00 & 0.00 & 0.00 \\
\$200K & 0.00 & 0.33 & 0.00 & 2.27 \\
\$500K & 3.20 & 4.68 & 3.32 & 8.40 \\
\$1000K & 7.72 & 9.81 & 7.93 & 14.27 \\
\bottomrule
\end{tabular}
\begin{tablenotes}
\small
\item CEV: consumption-equivalent variation (\% of lifetime consumption
agent would pay for annuity market access). Welfare model uses representative
SS income (\$18{,}500/yr, median quartile). Baseline: MWR=0.87, 2\% inflation,
$\psi=0.981$, $\psi_{\text{purchase}}=0.0163$. Group pricing: MWR=0.90.
Real annuity: 0\% inflation, MWR=0.87. Best feasible combines group pricing,
real annuity, no survival pessimism ($\psi=1.0$), and default architecture
($\psi_{\text{purchase}}=0$).
\end{tablenotes}
\end{table}


\begin{figure}[htbp]
\centering
\includegraphics[width=0.8\textwidth]{../figures/pdf/fig2_gamma_sensitivity.pdf}
\caption{Predicted ownership by coefficient of relative risk aversion ($\gamma$). All other parameters at baseline values ($\pi = \pInflation$, MWR $= \pMwrBaseline$, DFJ bequests). The shaded band marks observed ownership of \pctHRSLifetime--\pctHRSIannPooled.}
\label{fig:gamma}
\end{figure}

\begin{figure}[htbp]
\centering
\includegraphics[width=0.8\textwidth]{../figures/pdf/fig3_hazard_sensitivity.pdf}
\caption{Predicted ownership by hazard multiplier specification. Each specification varies the Good- and Poor-health multipliers while holding Fair at 1.0. Labels indicate the empirical source for each calibration. The shaded band marks observed ownership of \pctHRSLifetime--\pctHRSIannPooled.}
\label{fig:hazard}
\end{figure}

\begin{figure}[htbp]
\centering
\includegraphics[width=\textwidth]{../figures/pdf/fig4_policy_functions.pdf}
\caption{Optimal annuitization fraction ($\alpha^*$) at age 65 by initial wealth, for Good health ($H = 1$). Left panel: no bequest motive and no medical-expenditure risk, with pricing loads, inflation, and survival pessimism active. Right panel: full nine-channel structural model (DFJ bequests, medical-expenditure risk, Reichling--Smetters health-mortality correlation, survival pessimism, age-varying consumption needs, state-dependent utility, public-care aversion, pricing loads, and inflation). Both panels use MWR $= \pMwrBaseline$, $\pi = \pInflation$, and a uniform Social Security income at the sample mean.}
\label{fig:policy}
\end{figure}

\begin{figure}[htbp]
\centering
\includegraphics[width=0.8\textwidth]{../figures/pdf/fig5_cev_heatmap.pdf}
\caption{Consumption-equivalent variation (\%) from annuity market access, by initial wealth and health status. No-bequest specification ($\theta = 0$). Full model with medical costs, Reichling--Smetters correlation, MWR $= \pMwrBaseline$, and 2\% inflation.}
\label{fig:cev}
\end{figure}


\bibliography{bibliography}

\end{document}
